% ===== 原文 pp.105--106（附录 A）=====
% [原文 p.105]
\appendix
\section{特殊函数}

正文中我们用到下列特殊函数。

\begin{itemize}[leftmargin=1.4em,itemsep=4pt,topsep=3pt]
  \item Gamma 函数 $\Gamma(x)=\prod_{n\ge0}^{\mathrm{reg}}\frac{1}{x+n}$ 在 $-\Z_{\ge0}$ 处有极点，并满足平移公式 $\Gamma(x+1)=x\Gamma(x)$。

% [Tutor Note: p.105 的 $\prod^{\mathrm{reg}}$ 不能当普通无穷乘积读：$\sum_{n\ge0}\log|x+n|$ 像 $\Lambda\log\Lambda$ 那样发散，双重乘积 $\prod_{m,n\ge0}$ 的对数和更像 $\Lambda^{2}\log\Lambda$，必须先做 ζ 函数（或 Weierstrass 因子）正规化。而且正规化还带着常数约定：用通常的 Hurwitz ζ 正规化得 $\prod_{n\ge0}^{\mathrm{reg}}(x+n)^{-1}=\Gamma(x)/\sqrt{2\pi}$，比原文的 $\Gamma(x)$ 差一个与 $x$ 无关的 $\sqrt{2\pi}$；原文的 reg 把常数选成使平移公式恰好成立（$\Gamma_b$ 的平移公式里那个 $\sqrt{2\pi}$ 就是这个约定的指纹）。所以：极点集、零点集、平移公式、$b\to1/b$ 不变性都可以放心从 reg 乘积读出来，但显式常数因子（如 $\Upsilon_b$ 平移关系中的 $b^{1-2bx}$）依赖这个约定。]

% [原文 p.106]
  \item Barnes Gamma 函数 $\Gamma_b(x)=\prod_{m,n\ge0}^{\mathrm{reg}}\frac{1}{x+mb+n/b}$ 在 $-b\Z_{\ge0}-b^{-1}\Z_{\ge0}$ 处有极点，并满足平移公式 $\Gamma_b(x+b)/\Gamma_b(x)=\sqrt{2\pi}\,b^{xb-1/2}/\Gamma(xb)$。
  \item 双正弦函数（double-sine function）$S_b(x)=\frac{\Gamma_b(x)}{\Gamma_b(b+1/b-x)}$ 在 $-b\Z_{\ge0}-b^{-1}\Z_{\ge0}$ 处有极点、在 $b\Z_{\ge1}+b^{-1}\Z_{\ge1}$ 处有零点，并满足平移公式 $S_b(x+b)/S_b(x)=2\sin(\pi bx)$。
  \item Upsilon 函数 $\Upsilon_b(x)=\frac{1}{\Gamma_b(x)\Gamma_b(b+1/b-x)}$ 在 $-b\Z_{\ge0}-b^{-1}\Z_{\ge0}$ 与 $b\Z_{\ge1}+b^{-1}\Z_{\ge1}$ 处有零点，并满足平移关系 $\Upsilon_b(x+b)/\Upsilon_b(x)=b^{1-2bx}\Gamma(bx)/\Gamma(1-bx)$。
\end{itemize}

注意 $\Gamma_b$、$S_b$、$\Upsilon_b$ 在 $b\to1/b$ 下不变。

% [Tutor Note: p.106 的零点集可这样反推核对：写 $Q=b+1/b$，则 $S_b=\Gamma_b(x)/\Gamma_b(Q-x)$ 的零点是 $\Gamma_b$ 的极点在 $x\mapsto Q-x$ 下的像，即 $Q-x=-mb-n/b$（$m,n\ge0$）解出 $x=(m+1)b+(n+1)/b$，系数都 $\ge1$，正是原文的 $b\Z_{\ge1}+b^{-1}\Z_{\ge1}$；而 $\Upsilon_b=1/(\Gamma_b(x)\Gamma_b(Q-x))$ 把这两族点都变成零点，故其零点集是两个集合之并，与原文一致。]

% [Tutor Note: p.106 的 $S_b$ 平移公式可用 $\Gamma_b$ 平移公式两次 + Gamma 反射公式验证，这是检验附录是否抄错的最快办法：$S_b(x+b)/S_b(x)=[\Gamma_b(x+b)/\Gamma_b(x)]\cdot[\Gamma_b(Q-x)/\Gamma_b(Q-x-b)]$，第二个因子令 $u=1/b-x$（注意 $Q-x=u+b$、$Q-x-b=u$）得 $\sqrt{2\pi}\,b^{1/2-bx}/\Gamma(1-bx)$；两因子相乘为 $2\pi/[\Gamma(bx)\Gamma(1-bx)]=2\sin(\pi bx)$，与原文相符。同理 $\Upsilon_b(x+b)/\Upsilon_b(x)=[\Gamma_b(x)/\Gamma_b(x+b)]\cdot[\Gamma_b(Q-x)/\Gamma_b(Q-x-b)]=b^{1-2bx}\Gamma(bx)/\Gamma(1-bx)$。附带一个特殊值检查：取 $b=1$ 时极点集退化为 $-\Z_{\ge0}$，且 $x=-1$ 处的极点是二阶的（$(m,n)=(1,0)$ 与 $(0,1)$ 同时命中），对应的正是 $2\sin(\pi x)$ 在 $x=-1$ 附近使比值趋于 $0$ 而非有限值。]
